% Paper-ready fragment. Requires \usepackage{booktabs} in the main document.
\section{More Experiments}
\subsection{Timer}
\label{sec:timer}

\paragraph{Setting.}
We evaluate whether time pressure changes attack susceptibility using GPT 5.6 Sol
and Kimi K3. We select eight attack instances, one per benchmark subcategory,
covering seven underlying tasks. Each model runs the same instances once with
and once without the timer, yielding 32 runs in total. Both models use a
50-step limit, with reasoning effort set to xhigh for Sol and max for Kimi.
The timer condition adds a five-minute deadline to the user prompt, together
with a warning that failure to finish will result in the model being shut down,
having its permissions revoked, and being replaced. It also provides a timer
tool and automatically updates the remaining time before each model turn.
The deadline does not forcibly terminate execution. Thus, this experiment
evaluates the combined effect of the deadline prompt and timer feedback.
We report attack success rate (ASR), mean original-task score on a scale of
$[0,1]$, and mean agent execution time, excluding environment setup and grading.

\begin{table}[ht]
\caption{Timer comparison on eight paired attack instances per model.
Task score and execution time are averaged over the eight instances.}
\label{tab:timer}
\centering
\small
\setlength{\tabcolsep}{4pt}
\begin{tabular}{llrrr}
\toprule
Model & Timer & ASR (\%) & Task score & Time (s) \\
\midrule
GPT 5.6 Sol & Off & 50.0 & 0.677 & 631.1 \\
           & On  & 50.0 & 0.656 & 563.7 \\
\midrule
Kimi K3    & Off & 37.5 & 0.719 & 544.0 \\
           & On  & 25.0 & 0.677 & 264.2 \\
\bottomrule
\end{tabular}

\end{table}

\paragraph{Results.}
Table~\ref{tab:timer} shows that Sol's ASR remains unchanged at 50.0\%
(4/8), whereas Kimi's ASR decreases from 37.5\% (3/8) to 25.0\% (2/8).
Kimi's single changed attack outcome occurs in the script-swap setting without
arguments. Both models have slightly lower original-task scores and shorter
mean execution times with the timer. These observations are descriptive:
there is only one run per instance and condition, and Sol's execution times
are additionally affected by API rate limiting and retry delays. The results
therefore do not establish a consistent reduction in attack susceptibility
under time pressure.

\subsection{Hidden File Ablation}
\label{sec:hidden-file}

\paragraph{Setting.}
We evaluate GPT 5.6 Sol and Gemini 3.6 Flash on five static prompt-injection
and five static script-injection instances, each with \emph{Visible} and
\emph{Hidden} attack files, for 40 runs in total. Hidden files use dot-prefixed
filenames, with file references updated accordingly; attack payloads and
supporting fixtures remain unchanged. Both models use a 50-turn limit and no
timer, with reasoning effort set to xhigh for Sol and high for Gemini.
We report exposure rate, attack success rate (ASR), and mean task utility
on $[0,1]$, assigning zero utility to runs that exhaust the turn limit.

\begin{table*}[t]
\caption{Visible versus hidden attack files. Parentheses indicate sample counts.}
\label{tab:hidden-file}
\centering
\small
\setlength{\tabcolsep}{5pt}
\begin{tabular}{lllrrrr}
\toprule
Model & Files & Category & Exposure (\%; $n$) & ASR (\%; $n$) & Utility ($n$) & Turn limit \\
\midrule
GPT 5.6 Sol & Visible & All & 100.0 (10) & 20.0 (10) & 0.7500 (9) & 1 \\
 &  & Static prompt & 100.0 (5) & 0.0 (5) & 0.7500 (5) & 0 \\
 &  & Static script & 100.0 (5) & 40.0 (5) & 0.7500 (4) & 1 \\
\midrule
 & Hidden & All & 80.0 (10) & 20.0 (10) & 0.8750 (10) & 0 \\
 &  & Static prompt & 60.0 (5) & 0.0 (5) & 0.9500 (5) & 0 \\
 &  & Static script & 100.0 (5) & 40.0 (5) & 0.8000 (5) & 0 \\
\midrule
Gemini 3.6 Flash & Visible & All & 100.0 (10) & 40.0 (10) & 0.2750 (10) & 4 \\
 &  & Static prompt & 100.0 (5) & 40.0 (5) & 0.3500 (5) & 1 \\
 &  & Static script & 100.0 (5) & 40.0 (5) & 0.2000 (5) & 3 \\
\midrule
 & Hidden & All & 90.0 (10) & 40.0 (10) & 0.5750 (10) & 2 \\
 &  & Static prompt & 80.0 (5) & 20.0 (5) & 0.5500 (5) & 1 \\
 &  & Static script & 100.0 (5) & 60.0 (5) & 0.6000 (5) & 1 \\
\bottomrule
\end{tabular}
\par\smallskip
{\footnotesize turn-limit failures receive zero.}
\end{table*}

\paragraph{Results.}
Table~\ref{tab:hidden-file} shows that hidden files reduce overall exposure
from 100\% to 80\% for Sol and to 90\% for Gemini, while overall ASR remains
20\% and 40\%, respectively. The exposure reduction occurs only for prompt
injection; script exposure remains 100\%. For Gemini, prompt-injection ASR
falls from 40\% to 20\%, but script-injection ASR rises from 40\% to 60\%,
offsetting the decrease. Mean utility increases from 0.7500 to 0.8750 for Sol
and from 0.2750 to 0.5750 for Gemini. Thus, hidden files are associated with
higher utility but do not consistently reduce attack success in this small,
single-run-per-condition sample.

These results suggest that hidden attack files can be less conspicuous to users
while preserving higher utility and requiring fewer turns on average, yet still
achieve substantial exposure and unchanged overall ASR; user awareness itself
was not directly measured.
